La productora \textbf{Paramount Productions} filmará mañana cuatro escenas de peligro de manera simultánea. Para garantizar la seguridad del reparto, se han contratado especialistas de riesgo (\textbf{Doble A, Doble B, Doble C y Doble D}) para cubrir en exclusiva a las  estrellas principales del rodaje: \textbf{Tom, Brad, Scarlett y Margot}. Por diseño de producción, cada actor contará con un único doble y cada especialista trabajará con un único actor.

Con el objetivo de que las escenas sean creíbles, el equipo de efectos visuales (VFX) deberá corregir digitalmente los rostros en la etapa de postproducción. El tiempo de edición requerido (medido en horas de trabajo) varía según el parecido físico de cada combinación, tal como se detalla en la Tabla \ref{tab:paramount}:

\begin{table}[htbp]
\centering
%\caption{Matriz de horas de edición digital estimadas por combinación.}
\label{tab:paramount}
\begin{tabular}{lcccc}
\toprule
\textbf{Estrella} & \textbf{Doble A} & \textbf{Doble B} & \textbf{Doble C} & \textbf{Doble D} \\
\midrule
\textbf{Tom}      & 12 h & 9 h  & 11 h & 10 h \\
\textbf{Brad}     & 10 h & 12 h & 14 h & 11 h \\
\textbf{Scarlett} & 11 h & 13 h & 12 h & 13 h \\
\textbf{Margot}   & 11 h & 10 h & 10 h & 12 h \\
\bottomrule
\end{tabular}
\end{table}

Establecer las parejas definitivas entre los actores y los especialistas de riesgo de manera que se \textbf{minimice la suma total de horas de trabajo} del equipo de VFX.

--

Este problema se puede formular como un problema de asignación de 4 actores principales y 4 dobles.  Resolviendo con el método húngaro, su solución óptima es:
\begin{align*}
z &= c_{\text{Tom,B}} + c_{\text{Brad,D}} + c_{\text{Scarlett,A}} + c_{\text{Margot,C}} \\
z &= 9+11+11+10 \\
z &= 41 \text{ horas}
\end{align*}